\documentclass{article}
\usepackage[T1]{fontenc}
\usepackage{lmodern}
\usepackage{graphicx}
\usepackage{amsmath,amssymb}
\usepackage[a4paper,margin=2cm]{geometry}
\usepackage[absolute,overlay]{textpos}
\setlength{\TPHorizModule}{1mm}
\setlength{\TPVertModule}{1mm}
\usepackage[indonesian]{babel}
\usepackage{hyperref}
\setlength{\emergencystretch}{2em}
% unit: Halaman 18

\begin{document}

% === Halaman 18 (python/native) ===

18

• WSFP dipetakan menjadi ha*e. 

• Dalam Bahasa Inggris, kata yang mungkin untuk ha*e hanyalah have, hate, 

hale, dan haze 

• Dengan mencoba mengganti semua F di dalam cipherteks dengan v, t, l, dan z, 

maka huruf yang cocok adalah v sehingga WSFP dipetakan menjadi have

• Dengan mengganti F menjadi v pada  kriptogram EPYEPOPDZSZUFPO sehingga 

menjadi *e*e*e*tat*ve*, maka kata yang cocok untuk ini adalah 

representatives

\end{document}
